%%%%%
%
%% Dans ce fichier, les �ventuels accents sont cod�s
%% en ISO-8859-15 (aka latin9)
%% La ligne suivante donne trois exemples 
%% e accent aigu : �, e accent grave : �, e-dans-l'o : �
\section{Rotation horizontale appliqu�e � un vecteur vent pour passer dans
le rep�re du vent moyen}

\itemstyle{Op\'erateur~:} \operfmt{Rotat\_uv\_vent\_moy}
\index{Rotat\_uv\_vent\_moy}

%%%%%OBJET%%%%
\itemstyle{Objet~:}
Application au vecteur vent horizontal d'une rotation simple dans le rep�re dont le 1er axe est port� par le vent horizontal moyen.

%%%%%ARGUMENTS%%%%
\bigskip
%% Dans ce fichier, les �ventuels accents sont cod�s
%% en ISO-8859-15 (aka latin9)
%% La ligne suivante donne trois exemples 
%% e accent aigu : �, e accent grave : �, e-dans-l'o : �
%%%%%ENTREES%%%%
\begin{tabarg}
\multicolumn{3}{l}{\titretab {Entr\'ees}}\\
\hline
$U$ & Vecteur, \texttt{float} & 1�re composante du vent (unit� quelconque)\\
$V$ & Vecteur, \texttt{float} & 2�me composante du vent (m\^eme unit� que $U$)\\
$F_e$ & Coeff., \texttt{float} & Fr�quence d'�chantillonnage (Hz)\\
$T_{int}$ & Coeff., \texttt{float} & P�riode d'int�gration pour le calcul de la moyenne (s)\\
\end{tabarg}

\medskip
%%%%%SORTIES%%%%
\begin{tabarg}
\multicolumn{3}{l}{\titretab {Sorties}}\\
\hline
$U_r$ & Vecteur, \texttt{float} & Composante du vent dans l'axe du vent moyen (m\^eme unit� que $U$ et $V$)\\
$V_r$ & Vecteur, \texttt{float} & Composante du vent transverse (� gauche) � l'axe du vent moyen (m\^eme unit� que $U$ et $V$)\\
\end{tabarg}


%%%%%FORMULE%%%%
\itemstyle{Formule}

%% Formule de l'operateur : Rotat_uv_vent_moy
%% Dans ce fichier, les �ventuels accents sont cod�s
%% en ISO-8859-15 (aka latin9)
%% La ligne suivante donne trois exemples 
%% e accent aigu : �, e accent grave : �, e-dans-l'o : �




$$
\begin{cases}
U_r =     \cos \alpha .U  + \sin \alpha .V  \\
V_r = -\sin \alpha.U + \cos \alpha.V  \\
\end{cases}
$$

avec :\\
\begin{center}
$    \cos \alpha = \frac{U_m}{\sqrt{U_m^2 + V_m^2}} $ \\
$    \sin \alpha = \frac{V_m}{\sqrt{U_m^2 + V_m^2}} $ \\
\end{center}
\vspace{0.7cm}
$U_m$ et $V_m$ sont les moyennes de $U$ et $V$ sur la p�riode $T_{int}$.

%%%%%AUTEUR%%%%
\itemstyle{Auteur~:}
CNRM/GMEI/TRAMM

%%%%%REFERENCE%%%%
\itemstyle{R\'ef\'erences~:}
Pas n�cessaire.
